% Compiles on its own (pulls ../main.tex's preamble via subfiles) or as part of the book.
\documentclass[../main.tex]{subfiles}
\begin{document}
\noindent\fbox{\parbox{0.94\textwidth}{\textbf{Historical, quarantined
experiment.} These 16-epoch test-set numbers predate the corrected protocol and
must not support a claim. See \texttt{../../../EXPLANATION.md}.}}\par\medskip

\chapter{The Baselines and the Noise Floor}

Before a single fusion model exists, we need two things: the scores fusion must beat, and a
measure of how much a score can be trusted. This chapter produces both, and then does something
the rest of the field mostly skips --- it \emph{audits} the numbers before believing them.

\section{Why baselines come first}

The entire question of this project is: \emph{does combining two cameras beat using one?} That
question is meaningless without ``one'' as a reference. So the first experiments are deliberately
un-clever: train a detector on \textbf{thermal only}, train one on \textbf{RGB only}, and measure
them exactly the way every later model will be measured. These are the \textbf{baselines}. A fusion
model that cannot beat the better of them is not fusion --- it is a more expensive single camera,
and (as \texttt{apex} showed) that failure can hide behind a healthy-looking average.

\section{The noise floor, and why five seeds}

A single training run gives a single number, and that number is partly luck. The model starts from
random weights (from scratch --- Soham's directive), and different random starts land in slightly
different places. So we train each baseline \textbf{five times}, changing only the random seed, and
report the mean and the standard deviation.

That standard deviation is the \textbf{noise floor}: the amount a score wobbles for reasons that
have nothing to do with the model. It is the most important number for everything that follows,
because it sets the bar a real result must clear. The convention: a later model only counts as
genuinely better if it beats a baseline by more than \emph{twice} the noise floor
($|\Delta| > 2\sigma$). Anything smaller is indistinguishable from a lucky seed.

\begin{lesson}
Most published KAIST papers report single-run improvements of $0.3$--$1.0$ MR points with no error
bars at all. Our own noise floor (below) is $\approx 1$--$2$ points. That means many published
``improvements'' are inside the noise --- unfalsifiable. Running five seeds and quoting $\sigma$
makes this study more honest than much of the literature it cites. The discipline is cheap; the
credibility it buys is not.
\end{lesson}

\section{The results}

Five seeds each, from scratch, at the locked recipe (lr $=0.01$, 16 epochs), evaluated on the
improved-test annotations:

\begin{center}
\begin{tabular}{lccc}
\toprule
 & All & Day & Night \\
\midrule
\textbf{thermal-only} & $36.1\pm1.3$ & $43.9\pm1.4$ & $\mathbf{20.6\pm1.2}$ \\
\textbf{RGB-only}     & $52.5\pm1.5$ & $46.3\pm1.8$ & $\mathbf{65.9\pm0.9}$ \\
\bottomrule
\end{tabular}
\end{center}

\begin{center}
\includegraphics[width=0.78\textwidth]{baselines_thermal_vs_rgb.png}
\end{center}

The headline is \textbf{Night}: RGB misses $\approx 66\%$ of pedestrians, thermal $\approx 21\%$ ---
a $45$-point chasm. At night the RGB image is nearly black; thermal sees body heat glowing against
the cool dark. This is the multispectral premise made concrete, and it is the gap fusion exists to
exploit. Thermal is in fact the stronger camera \emph{everywhere} here (even in daytime), so RGB
has no clear home turf --- which makes fusion's job specific: add RGB's daytime detail on top of
thermal's dominance \emph{without} surrendering thermal's night advantage.

\section{Auditing the numbers before trusting them}

A clean-looking result is exactly when to be suspicious. Five checks, each able to expose a
different bug:

\begin{enumerate}
  \item \textbf{Internal consistency.} The ``All'' score should approximate the day/night scores
        blended by their frame counts ($1455$ day, $797$ night). Thermal: blend $35.6$ vs actual
        $36.1$; RGB: blend $53.2$ vs actual $52.5$. Both agree within a point. \emph{This is the
        check that catches a broken evaluator, a wrong frame count, or a swapped day/night label ---
        and it passes.}
  \item \textbf{Cameras not swapped.} The worst silent bug would feed the ``thermal'' model RGB
        images. But thermal is better at night and RGB is better in day --- both physically correct.
        A swap would make ``thermal'' blind at night. Ruled out.
  \item \textbf{Night ordering vs physics.} The $45$-point night gap is the textbook multispectral
        result; a broken pipeline gives $\sim100\%$ (detects nothing) or noise, not a clean
        physics-consistent split.
  \item \textbf{RGB night vs the published floor.} The literature's RGB-only night is $\approx 64\%$;
        ours is $65.9\%$. At night RGB hits a floor --- you cannot detect what the camera cannot see
        --- so it lands there almost regardless of architecture. Independent agreement.
  \item \textbf{Absolute levels.} Thermal ($36$) sits between 2016 Faster-RCNN-thermal ($\sim47$)
        and modern pretrained thermal ($\sim19$); RGB ($52$) is weaker than pretrained
        Faster-RCNN-RGB, as from-scratch should be. Believable, not suspicious.
\end{enumerate}

\begin{lesson}
One result deviates from the textbook pattern: in the literature RGB usually beats thermal in
\emph{daytime}, but here thermal marginally wins day too ($43.9$ vs $46.3$). This is not a bug --- it
is a from-scratch artifact. Thermal images are simpler (grayscale-like heat blobs) and therefore
easier to learn from a random start, while RGB (rich colour and texture) underperforms its usual
potential without a pretrained backbone. Pretrained models do not show this because ImageNet already
taught the RGB backbone to see. The margin is $\approx 2\sigma$ --- marginal, well-explained, worth
a footnote in the report, not worth worrying about. The habit to keep: when a number defies
expectation, find the mechanism before you either trust it or ``fix'' it.
\end{lesson}

\section{What this licenses}

With trustworthy baselines and a quantified noise floor, the study has its yardstick. Every fusion
model from here reports MR$^{-2}$ All/Day/Night against \emph{these} numbers, and a claim of
improvement is only entered in the ledger if it clears $2\sigma$. The next chapter builds the first
model that sees both cameras at once.

\end{document}
